% GENERATED FILE. Edit canonical lessons, then run npm run generate:books.
% canonical-source-hash: fnv1a64:01929c4744297677
% canonical-lessons: BN-W02-lal-read, BN-W03-kalo-read, BN-W03-shada-read, BN-W04-ii-matra, BN-W04-nil-read, BN-W02-tin-read, BN-W04-ek-read

\chapter{The Last Sign, and Six Words Back}
\label{ch:bn-last-sign-six-words}
\hlchaptermodality{24}

\begin{chapteropening}
\textbf{By the end of this chapter:} \emph{I can read the colours and two of the numerals off the page, write the long-i sign that completes the script this book teaches, and say why every one of these words reached my ear before it reached my eye.}

One read as a vowel opening a word with no consonant to carry it, which finally makes visible the claim the numbers chapter could only assert; then twenty-six pieces and nineteen words counted back, all of them heard before they were seen.
\end{chapteropening}

\section[lāl]{\bn{লাল} — red, and the day-word that runs both ways}

\label{lesson:BN-W02-lal-read}



\begin{quote}
Write \textbf{\bn{ল}} once, then say where the tongue sits for it.
\end{quote}

\subsection*{Script}
\begin{quote}
\textbf{\bn{ল}} + \textbf{\bn{া}} + \textbf{\bn{ল}} $\to$ \textbf{\bn{লাল}}
\end{quote}

The same letter twice, with the mātrā between them. The first \textbf{\bn{ল}} takes the sign and becomes \emph{lā}; the second keeps its own inherent vowel, which at the end of a word is barely voiced at all. So the word comes out \textbf{lāl} — one syllable to an English ear, the colour \textbf{red}, the Persian gem-word from the chapter on colours.

Swap the first shape for \textbf{\bn{ক}} and you get another word already in your mouth:

\begin{quote}
\textbf{\bn{ক}} + \textbf{\bn{া}} + \textbf{\bn{ল}} $\to$ \textbf{\bn{কাল}}
\end{quote}

\emph{kāl} — the word that means \textbf{tomorrow} and \textbf{yesterday} both, and which the farewell chapter used in \emph{kāl dækhā hôbe}. Bengali does not think the two days need separate names; the distance from today is what the word records, and the direction comes from the verb.

Two words, three pieces each, one shape different between them. Reading is starting to be recombination rather than recognition.

\subsection*{Your turn}
\begin{itemize}
\raggedright
  \item \emph{Read it:} \bn{লাল}
  \item \emph{Read it:} \bn{কাল}
  \item \emph{Read it:} \bn{আমি}
  \item \emph{Write it:} \bn{লাল}
\end{itemize}

\begin{tcolorbox}[breakable,colback=teal!4,colframe=teal!35!black,title={Before you move on}]
Why is the second \textbf{\bn{ল}} in \textbf{\bn{লাল}} not said \emph{lô}? (Its inherent vowel is \textbf{barely voiced at the end of a word}.) And which two days does \textbf{\bn{কাল}} cover? (\textbf{Yesterday and tomorrow} — both.)

Source: \href{https://www.unicode.org/charts/PDF/U0980.pdf}{Unicode Bengali chart}.
\end{tcolorbox}
\section[kālo]{\bn{কালো} — black, read}

\label{lesson:BN-W03-kalo-read}



\begin{quote}
Write \textbf{\bn{ো}}, and name the two signs it is made of.
\end{quote}

\subsection*{Script}
\begin{quote}
\textbf{\bn{ক}} + \textbf{\bn{া}} + \textbf{\bn{ল}} + \textbf{\bn{ো}} $\to$ \textbf{\bn{কালো}}
\end{quote}

\begin{itemize}
\raggedright
  \item \textbf{\bn{ক}} with \textbf{\bn{া}} after it — \emph{kā}
  \item \textbf{\bn{ল}} with \textbf{\bn{ো}} around it — \emph{lo}
\end{itemize}

\textbf{kālo} — \emph{black}, the colour that is the same word as Sanskrit \emph{kāla}, \textquotedblleft{}time,\textquotedblright{} and stands behind both Kālī and Yama.

This is the first word in the book with \textbf{two different vowel signs on two consecutive letters}, and it is the moment the eye has to stop scanning and start parsing. There are three marks after the first \textbf{\bn{ক}} and only two of them belong together. Work it as the last lesson said: find the consonant, then look around it.

Set it beside the red from the last chapter and the two are nearly twins on the page:

\begin{quote}
\textbf{\bn{লাল}} and \textbf{\bn{কালো}}
\end{quote}

Both open with a consonant plus \textbf{\bn{া}}; both close on \textbf{\bn{ল}}. What separates them is which letter the \textbf{\bn{ল}} is and what is hanging on it.

\subsection*{Your turn}
\begin{itemize}
\raggedright
  \item \emph{Read it:} \bn{কালো}
  \item \emph{Read it:} \bn{লাল}
  \item \emph{Read it:} \bn{কেমন}
  \item \emph{Write it:} \bn{কালো}
\end{itemize}

\begin{tcolorbox}[breakable,colback=teal!4,colframe=teal!35!black,title={Before you move on}]
Which letter in \textbf{\bn{কালো}} is wearing the o-sign? (The \textbf{\bn{ল}} — the second one.) And what else does the same Sanskrit word behind it mean? (\textbf{Time}.)

Source: \href{https://www.unicode.org/charts/PDF/U0980.pdf}{Unicode Bengali chart}.
\end{tcolorbox}
\section[shādā]{\bn{সাদা} — white, read}

\label{lesson:BN-W03-shada-read}



\begin{quote}
Write \textbf{\bn{দ}}, and say what separates it from \textbf{\bn{ত}}.
\end{quote}

\subsection*{Script}
\begin{quote}
\textbf{\bn{স}} + \textbf{\bn{া}} + \textbf{\bn{দ}} + \textbf{\bn{া}} $\to$ \textbf{\bn{সাদা}}
\end{quote}

Two consonants, each wearing the same mātrā. As plain a spelling as this book has shown.

\textbf{shādā} — \emph{white}. And notice the first sound: the letter is the \emph{s}-letter, and it is said \textbf{sh}. That was the news of the greeting back in the first script chapter, and it has not changed. In Bengali this letter is \emph{sh} far more often than it is \emph{s}, which is the second leg of the track's sound-fingerprint.

Three colours are now readable:

\begin{quote}
\textbf{\bn{লাল}} \emph{lāl} — red \textbf{\bn{কালো}} \emph{kālo} — black \textbf{\bn{সাদা}} \emph{shādā} — white
\end{quote}

White is the one of the three where Bengali and Hindi went different ways: Hindi took a Persian loan for it and Bengali kept its own word, where for red both took the same Persian gem-name. A page of Bengali colour words is a small map of who borrowed what.

\subsection*{Your turn}
\begin{itemize}
\raggedright
  \item \emph{Read it:} \bn{সাদা}
  \item \emph{Read it:} \bn{লাল} \bn{কালো} \bn{সাদা}
  \item \emph{Write it:} \bn{সাদা}
\end{itemize}

\begin{tcolorbox}[breakable,colback=teal!4,colframe=teal!35!black,title={Before you move on}]
The first letter of \textbf{\bn{সাদা}} is written \emph{s}. How is it said? (\emph{\textbf{sh}}.) And which of the three colours is the one Bengali kept where Hindi borrowed? (\textbf{\bn{সাদা}} — white.)

Source: \href{https://www.unicode.org/charts/PDF/U0980.pdf}{Unicode Bengali chart}.
\end{tcolorbox}
\section[-i (long)]{\bn{ী} — the second i-sign}

\label{lesson:BN-W04-ii-matra}



\begin{quote}
Read \textbf{\bn{জল}}, then write \textbf{\bn{ি}} and say which side it goes on.
\end{quote}

\subsection*{Script}
\begin{quote}
\bn{ী}
\end{quote}

A second sign for \textbf{i} — and in modern Bengali speech it sounds \textbf{exactly like the first one}.

Sanskrit distinguished a short \emph{i} from a long \emph{ī}, and held words apart by that length alone. Bengali gave the distinction up. The two signs survive because the spellings survive: a word inherited from Sanskrit keeps whichever sign it always had, and the reader learns it word by word rather than by ear.

That is worth stating plainly rather than apologising for. \textbf{A script records the history of a language as well as its present}, and the cost of that is a handful of spellings a listener could never have guessed. English readers know the arrangement well — the \emph{k} in \emph{know}, the \emph{gh} in \emph{night}.

The two signs sit on opposite sides:

\begin{quote}
\textbf{\bn{ি}} goes \textbf{before} its consonant \textbf{\bn{ী}} goes \textbf{after} it, like \textbf{\bn{া}}
\end{quote}

So they are hard to confuse on the page even while being impossible to tell apart in the ear.

\begin{quote}
\textbf{\bn{ন}} + \textbf{\bn{ী}} $\to$ \textbf{\bn{নী}}
\end{quote}

\subsection*{Your turn}
\begin{itemize}
\raggedright
  \item \emph{Write it:} \bn{ী}
  \item \emph{Write it:} \bn{নী}
  \item \emph{Read it:} \bn{জল}
\end{itemize}

\begin{tcolorbox}[breakable,colback=teal!4,colframe=teal!35!black,title={Before you move on}]
Do \textbf{\bn{ি}} and \textbf{\bn{ী}} sound different in Bengali today? (\textbf{No}.) Did they in Sanskrit? (\textbf{Yes} — short and long.) And which side does \textbf{\bn{ী}} go on? (\textbf{After} its consonant.)

Source: \href{https://www.unicode.org/charts/PDF/U0980.pdf}{Unicode Bengali chart}.
\end{tcolorbox}
\section[nīl]{\bn{নীল} — blue, read}

\label{lesson:BN-W04-nil-read}



\begin{quote}
Write \textbf{\bn{ী}}, and say which side of its consonant it goes on.
\end{quote}

\subsection*{Script}
\begin{quote}
\textbf{\bn{ন}} + \textbf{\bn{ী}} + \textbf{\bn{ল}} $\to$ \textbf{\bn{নীল}}
\end{quote}

\textbf{nīl} — \emph{blue}, the colour from Sanskrit \emph{nīla}, the dye that gave English \emph{indigo} its trade and gave the Bengal countryside its indigo plantations.

The sign in the middle is the long one, and that is not a choice a speaker could make. Sanskrit \emph{nīla} had a long \emph{ī}, so Bengali writes \textbf{\bn{ী}} — even though a Bengali speaker says it no differently from the short sign. The spelling is inherited, and the previous lesson's point is now something you have to act on rather than agree with.

Four colours are now readable:

\begin{quote}
\textbf{\bn{লাল}} — red \textbf{\bn{কালো}} — black \textbf{\bn{সাদা}} — white \textbf{\bn{নীল}} — blue
\end{quote}

Of these, blue is the one whose spelling carries a fact about Sanskrit rather than about Bengali. Red is a Persian gem-name; black is the word for time; white is the homegrown one Hindi replaced. Four words, four different histories, and the page shows one of them.

\subsection*{Your turn}
\begin{itemize}
\raggedright
  \item \emph{Read it:} \bn{নীল}
  \item \emph{Read it:} \bn{লাল} \bn{কালো} \bn{সাদা} \bn{নীল}
  \item \emph{Write it:} \bn{নীল}
\end{itemize}

\begin{tcolorbox}[breakable,colback=teal!4,colframe=teal!35!black,title={Before you move on}]
Why does \textbf{\bn{নীল}} take the long i-sign? (Because \textbf{Sanskrit \emph{nīla} had a long i}.) Can a speaker hear the difference? (\textbf{No}.) And how many colours can you read now? (\textbf{Four}.)

Source: \href{https://www.unicode.org/charts/PDF/U0980.pdf}{Unicode Bengali chart}.
\end{tcolorbox}
\section[tin]{\bn{তিন} — three, read}

\label{lesson:BN-W02-tin-read}



\begin{quote}
Write \textbf{\bn{ত}}, then write \textbf{\bn{ি}}, and say which side of a consonant the second one goes on.
\end{quote}

\subsection*{Script}
\begin{quote}
\textbf{\bn{ত}} + \textbf{\bn{ি}} + \textbf{\bn{ন}} $\to$ \textbf{\bn{তিন}}
\end{quote}

Read it the way the mouth needs it, not the way the eye meets it:

\begin{itemize}
\raggedright
  \item \textbf{\bn{ত}} — the dental \emph{t}, tongue on the teeth
  \item \textbf{\bn{ি}} — leaning on that \textbf{\bn{ত}} from the left, heard after it
  \item \textbf{\bn{ন}} — \emph{nô}, its inherent vowel almost silent at the end
\end{itemize}

\textbf{tin} — \emph{three}, counted since the numbers chapter.

The hook in front of \textbf{\bn{ত}} is the same trap as \textbf{\bn{আমি}}, and it will keep being the same trap until the hand has written it twenty times. On the page: hook, t, n. In the mouth: t, i, n.

This word is also about as old as words get. Sanskrit \emph{tri}, Latin \emph{trēs}, Greek \emph{treîs}, English \emph{three} — the numbers chapter set that out, and the point here is only that you can now meet it on a page in Kolkata instead of in a table.

\subsection*{Your turn}
\begin{itemize}
\raggedright
  \item \emph{Read it:} \bn{তিন}
  \item \emph{Read it:} \bn{লাল}
  \item \emph{Write it:} \bn{তিন}
\end{itemize}

\begin{tcolorbox}[breakable,colback=teal!4,colframe=teal!35!black,title={Before you move on}]
In \textbf{\bn{তিন}}, is the hook read before or after the \textbf{\bn{ত}}? (\textbf{After} — it is only written before.) And what does the word mean? (\textbf{Three}.)

Source: \href{https://www.unicode.org/charts/PDF/U0980.pdf}{Unicode Bengali chart}.
\end{tcolorbox}
\section[ek]{\bn{এক} — one, and the words this reader can now take off a page}

\label{lesson:BN-W04-ek-read}



\begin{quote}
Write \textbf{\bn{এ}} and \textbf{\bn{ই}}, and say what decides when a vowel gets its own letter.
\end{quote}

\subsection*{Script}
\begin{quote}
\textbf{\bn{এ}} + \textbf{\bn{ক}} $\to$ \textbf{\bn{এক}}
\end{quote}

\textbf{ek} — \emph{one}. Two pieces, and the first of them is a vowel standing on its own because a word cannot open on a sign.

The numbers chapter made a remark that could not be checked at the time: it said those five numerals were no good for showing off the inherent \textbf{ô}, because \textbf{\bn{এক}} opens on an independent vowel. Here is what that meant. There is no consonant at the front of this word to carry a vowel, so there is no inherent vowel to hear. The claim is now something you can see rather than take on trust.

And \textbf{\bn{তিন}}, read two lessons ago, is on the page too, so two of the five numerals are now shapes rather than sounds.

Count the whole holding. Twenty-six pieces, gathered nine, four, four, eight and one across five script chapters. And with them, nineteen words:

\begin{itemize}
\raggedright
  \item \textbf{\bn{নমস্কার}} \emph{the greeting} · \textbf{\bn{নাম}} \emph{name} · \textbf{\bn{তুমি}} \emph{you} · \textbf{\bn{আমি}} \emph{I} · \textbf{\bn{কেমন}} \emph{how}
  \item \textbf{\bn{আবার}} \emph{again} · \textbf{\bn{ভালো}} \emph{good} · \textbf{\bn{চা}} \emph{tea} · \textbf{\bn{জল}} \emph{water} · \textbf{\bn{দুধ}} \emph{milk}
  \item \textbf{\bn{পরিবার}} \emph{family} · \textbf{\bn{ভাই}} \emph{brother} · \textbf{\bn{মুখ}} \emph{mouth} · \textbf{\bn{লাল}} \emph{red} · \textbf{\bn{কালো}} \emph{black}
  \item \textbf{\bn{সাদা}} \emph{white} · \textbf{\bn{নীল}} \emph{blue} · \textbf{\bn{তিন}} \emph{three} · \textbf{\bn{এক}} \emph{one}
\end{itemize}

with \textbf{\bn{নাক}}, \textbf{\bn{কি}}, \textbf{\bn{আসি}}, \textbf{\bn{আমার}} and \textbf{\bn{কাল}} riding along free.

Every one of them was in the mouth before it was on the page. That order was the whole design: hear a word, use it, and only then be shown what it looks like. A reader who had been handed the alphabet first would have spent twenty-six lessons decoding shapes that meant nothing yet. Every script chapter in this book sits between two chapters of speech for that reason.

\subsection*{Your turn}
\begin{itemize}
\raggedright
  \item \emph{Read it:} \bn{মুখ} \bn{লাল} \bn{কালো} \bn{সাদা} \bn{নীল}
  \item \emph{Read it:} \bn{এক} \bn{তিন}
  \item \emph{Write it:} \bn{এক}
  \item \emph{Write it:} \bn{এ} \bn{ই} \bn{আ} \bn{ী} \bn{ি} \bn{া} \bn{ে} \bn{ো}
\end{itemize}

\begin{tcolorbox}[breakable,colback=teal!4,colframe=teal!35!black,title={Before you move on}]
Why is there no inherent vowel to hear at the start of \textbf{\bn{এক}}? (It opens on an \textbf{independent vowel}, not a consonant.) How many pieces do you hold now? (\textbf{Twenty-six}.) How many words? (\textbf{Nineteen}.) And which came first for every one of them — the sound or the shape? (\textbf{The sound}.)

Source: \href{https://www.unicode.org/charts/PDF/U0980.pdf}{Unicode Bengali chart}.
\end{tcolorbox}
